% Part: sets-functions-relations
% Chapter: functions
% Section: functions-relations

\documentclass[../../../include/open-logic-section]{subfiles}


\begin{document}

\olfileid[ar]{sfr}{fun}{rel}

\olsection{الدوال بوصفها علاقات}

\begin{explain}
الدالة التي تقابل !!{element}s~${\OLId{latin-looped}{A}}$ بـ!!{element}s~${\OLId{latin-looped}{B}}$ تعيّن علاقة
بين ${\OLId{latin-looped}{A}}$ و~${\OLId{latin-looped}{B}}$: وهي التي تقوم بين ${\OLId{latin-ordinary}{x}}$ و${\OLId{latin-ordinary}{y}}$ إذا وفقط إذا كان
${\OLId{latin-ordinary}{f}}({\OLId{latin-ordinary}{x}}) = {\OLId{latin-ordinary}{y}}$. وقد نريد، اختزالًا لأصول البناء الرياضي مثلًا،
\emph{مماهاة}~${\OLId{latin-ordinary}{f}}$ بهذه العلاقة، أي بمجموعة الأزواج.
فأي العلاقات يجوز أن تؤخذ دوال على هذا الوجه؟
\end{explain}

\begin{defn}[الرسم البياني لدالة]
للدالة ${\OLId{latin-ordinary}{f}}\colon {\OLId{latin-looped}{A}} \to {\OLId{latin-looped}{B}}$ \emph{رسم بياني} للدالة~${\OLId{latin-ordinary}{f}}$، وهو العلاقة
${\OLId{latin-looped}{R}}_{\OLId{latin-ordinary}{f}} \subseteq {\OLId{latin-looped}{A}} \times {\OLId{latin-looped}{B}}$ ذات التعريف
\[
{\OLId{latin-looped}{R}}_{\OLId{latin-ordinary}{f}} = \Setabs{\tuple{{\OLId{latin-ordinary}{x}},{\OLId{latin-ordinary}{y}}}}{{\OLId{latin-ordinary}{f}}({\OLId{latin-ordinary}{x}}) = {\OLId{latin-ordinary}{y}}}.
\]
\end{defn}

\begin{explain}
تقضي الامتدادية بوحدة رسم الدالة البياني. بل إن امتدادية المجموعات
تسوّغ مباشرة امتدادية الدوال التي استعملناها ضمنًا: متى اتحد
مجال ${\OLId{latin-ordinary}{f}}$ و~${\OLId{latin-ordinary}{g}}$ ومجالهما المقابل، واتفقت قيمهما كلها، كانتا واحدة.

وبالعكس، كل علاقة تستوفي شرط «الدالية» هي رسم بياني لدالة.
\end{explain}

\begin{prop}\ollabel{prop:graph-function}
افرض أن ${\OLId{latin-looped}{R}} \subseteq {\OLId{latin-looped}{A}} \times {\OLId{latin-looped}{B}}$ تحقق الشرطين:
\begin{enumerate}
\item اجتماع $Rxy$ و$Rxz$ يقتضي ${\OLId{latin-ordinary}{y}} = {\OLId{latin-ordinary}{z}}$؛ و
\item يوجد لكل ${\OLId{latin-ordinary}{x}} \in {\OLId{latin-looped}{A}}$ عنصر ${\OLId{latin-ordinary}{y}} \in {\OLId{latin-looped}{B}}$ يحقق $\tuple{{\OLId{latin-ordinary}{x}},
{\OLId{latin-ordinary}{y}}} \in {\OLId{latin-looped}{R}}$.
\end{enumerate}
فتكون ${\OLId{latin-looped}{R}}$ رسم الدالة ${\OLId{latin-ordinary}{f}}\colon {\OLId{latin-looped}{A}} \to {\OLId{latin-looped}{B}}$ التي نضعها بالشرط:
${\OLId{latin-ordinary}{f}}({\OLId{latin-ordinary}{x}}) = {\OLId{latin-ordinary}{y}}$ إذا وفقط إذا كان $Rxy$.
\end{prop}

\begin{proof}
إذا وجد ${\OLId{latin-ordinary}{y}}$ يحقق $Rxy$ امتنع أن يوجد ${\OLId{latin-ordinary}{z}} \neq
{\OLId{latin-ordinary}{y}}$ يحقق $Rxz$، وإلا خولف شرط~${\OLId{latin-looped}{R}}$. فمتى وجد ${\OLId{latin-ordinary}{y}}$ مع $Rxy$ كان
هذا ${\OLId{latin-ordinary}{y}}$ واحدًا؛ ومع ضمان الوجود تكون ${\OLId{latin-ordinary}{f}}$ سليمة التعريف.
ومن تعريفها يظهر ${\OLId{latin-looped}{R}}_{\OLId{latin-ordinary}{f}} = {\OLId{latin-looped}{R}}$.
\end{proof}

\begin{explain}
فكل دالة ${\OLId{latin-ordinary}{f}}\colon {\OLId{latin-looped}{A}} \to {\OLId{latin-looped}{B}}$ لها رسم، أي علاقة بين ${\OLId{latin-looped}{A}}$ و~${\OLId{latin-looped}{B}}$
محتواة في~${\OLId{latin-looped}{A}} \times {\OLId{latin-looped}{B}}$ يضبطها ${\OLId{latin-ordinary}{f}}({\OLId{latin-ordinary}{x}}) = {\OLId{latin-ordinary}{y}}$. وكل علاقة~${\OLId{latin-looped}{R}}
\subseteq {\OLId{latin-looped}{A}} \times {\OLId{latin-looped}{B}}$ تحقق شرطي \olref{prop:graph-function} رسم
لدالة~${\OLId{latin-ordinary}{f}} \colon {\OLId{latin-looped}{A}} \to {\OLId{latin-looped}{B}}$. ولشدة هذه الصلة يجوز أن نأخذ الدالة
رسمها نفسه، أي علاقة مخصوصة أو مجموعة مخصوصة من الصفوف المرتبة.
\oliflabeldef{sfr:rel:ref:sec}{والمقصود بالمماهاة هنا ما بُيّن في
\olref[sfr][rel][ref]{sec}: ليست دعوى ميتافيزيقية في ماهية الدوال،
بل إجازة نافعة لـ\emph{معاملتها} مجموعات معينة. ومن نفع ذلك أن }{}
العمليات التي أجريناها على العلاقات يمكن إجراء نظائرها على الدوال
(انظر \olref[sfr][rel][ops]{sec})؛ ومنها ما يأتي.
\end{explain}

\begin{defn}\ollabel{defn:funimage}
خذ دالة ${\OLId{latin-ordinary}{f}} \colon {\OLId{latin-looped}{A}} \to {\OLId{latin-looped}{B}}$ وجزئية ${\OLId{latin-looped}{C}}\subseteq {\OLId{latin-looped}{A}}$.

\emph{تقييد}~${\OLId{latin-ordinary}{f}}$ على~${\OLId{latin-looped}{C}}$ هو الدالة~$\funrestrictionto{{\OLId{latin-ordinary}{f}}}{{\OLId{latin-looped}{C}}}\colon {\OLId{latin-looped}{C}} \to {\OLId{latin-looped}{B}}$
التي تحقق $(\funrestrictionto{{\OLId{latin-ordinary}{f}}}{{\OLId{latin-looped}{C}}})({\OLId{latin-ordinary}{x}}) = {\OLId{latin-ordinary}{f}}({\OLId{latin-ordinary}{x}})$ لكل ${\OLId{latin-ordinary}{x}} \in {\OLId{latin-looped}{C}}$؛ أي
$\funrestrictionto{{\OLId{latin-ordinary}{f}}}{{\OLId{latin-looped}{C}}} = \Setabs{\tuple{{\OLId{latin-ordinary}{x}}, {\OLId{latin-ordinary}{y}}} \in {\OLId{latin-looped}{R}}_{\OLId{latin-ordinary}{f}}}{{\OLId{latin-ordinary}{x}} \in
{\OLId{latin-looped}{C}}}$.

و\emph{تطبيق}~${\OLId{latin-ordinary}{f}}$ على~${\OLId{latin-looped}{C}}$ تعريفه $\funimage{{\OLId{latin-ordinary}{f}}}{{\OLId{latin-looped}{C}}} =
\Setabs{{\OLId{latin-ordinary}{f}}({\OLId{latin-ordinary}{x}})}{{\OLId{latin-ordinary}{x}} \in {\OLId{latin-looped}{C}}}$؛ ويسمى أيضًا \emph{صورة}~${\OLId{latin-looped}{C}}$ تحت~${\OLId{latin-ordinary}{f}}$.
\end{defn}

\begin{explain}
ومن التعريفين ينتج $\ran{{\OLId{latin-ordinary}{f}}} =
\funimage{{\OLId{latin-ordinary}{f}}}{\dom{{\OLId{latin-ordinary}{f}}}}$ لكل دالة~${\OLId{latin-ordinary}{f}}$.
\oliflabeldef{sfr:rel:ops:sec}{وهذا نظير لما تقتضيه تعريفات
\olref[sfr][rel][ops]{sec} ومماهاة الدوال بالعلاقات، إلا أن تقييد
الدالة هنا إنما يقصر المدخل على~${\OLId{latin-looped}{C}}$ ولا يقصر المخرج على~${\OLId{latin-looped}{C}}$.
وأما المعكوسات والضروب النسبية فيحتاجان إلى زيادة تفصيل،
وسيأتي في \olref[inv]{sec} و\olref[cmp]{sec}.}{}
\end{explain}

\begin{explain}
\textbf{تنبيهان على لفظ الأصل.} سمى الأصل رسم~${\OLId{latin-ordinary}{f}}$ «علاقة على
${\OLId{latin-looped}{A}}\times {\OLId{latin-looped}{B}}$»، مع أن اصطلاح الكتاب يجعل العلاقة على~${\OLId{latin-looped}{U}}$ جزءًا من
${\OLId{latin-looped}{U}}^{\OLMathNumeral{2}}$؛ فقلنا علاقة بين ${\OLId{latin-looped}{A}}$ و~${\OLId{latin-looped}{B}}$ محتواة في~${\OLId{latin-looped}{A}}\times {\OLId{latin-looped}{B}}$. ثم إن
تقييد الدالة على~${\OLId{latin-looped}{C}}$ يخص الإحداثي الأول، فلا يطابق تقييد العلاقة
المتجانسة بالتقاطع مع~${\OLId{latin-looped}{C}}^{\OLMathNumeral{2}}$ إلا على جهة النظير.
\end{explain}

\end{document}
